\section{问题一：常热物性下的短时温湿分布}
\subsection{指定位置与时刻的计算结果}
表\ref{tab:paper1}和表\ref{tab:paper2}列出100、300、600、900、1200、1500、1800秒时，半径0、0.5、1.0、1.5、2.0厘米处的温度与干基含水率。结果由最终640区间网格的未舍入数值解采样得到，完整逐秒数据保存在工作簿1中。

\input{tables/table1.tex}
\begin{table}[htbp]
\centering\small
\caption{问题一指定时刻的干基含水率（对应题目表2，单位：kg/kg）}\label{tab:paper2}
\begin{tabular}{rrrrrr}
\toprule
时间/s & 0 cm & 0.5 cm & 1.0 cm & 1.5 cm & 2.0 cm \\
\midrule
100 & 2.5500 & 2.5500 & 2.5500 & 2.5500 & 2.2470 \\
300 & 2.5500 & 2.5500 & 2.5500 & 2.5492 & 2.0508 \\
600 & 2.5500 & 2.5500 & 2.5500 & 2.5353 & 1.8770 \\
900 & 2.5500 & 2.5500 & 2.5497 & 2.5045 & 1.7547 \\
1200 & 2.5500 & 2.5500 & 2.5482 & 2.4646 & 1.6586 \\
1500 & 2.5500 & 2.5499 & 2.5445 & 2.4206 & 1.5787 \\
1800 & 2.5500 & 2.5497 & 2.5383 & 2.3755 & 1.5102 \\
\bottomrule
\end{tabular}
\end{table}



0.50小时，中心和表面温度分别为33.5753、36.7856摄氏度，相差3.2102摄氏度；中心含水率为2.5500，接近初值，表面已降至1.5102。热量已影响整个截面，明显失水主要发生在近表面区域。

图\ref{fig:q1-profiles}给出100、600和1800秒的径向剖面。温度在中心附近平坦，向表面升高；含水率则在表面最低，向中心逐渐接近初值。两组剖面反映了热量传入与水分排出的不同进程。

\begin{figure}[htbp]\centering
\includegraphics[width=\linewidth]{figures/result_q1_profiles.pdf}
\caption{问题一在100、600和1800秒的径向温湿剖面，同一时刻采用相同颜色与线型。}\label{fig:q1-profiles}
\end{figure}

\subsection{传递尺度及边界阻力解释}
问题一的热扩散率为$k/(\rho_h c_p)=1.6886\times10^{-7}$平方米每秒，初始水分扩散系数为$4.9377\times10^{-9}$平方米每秒，前者约为后者的34倍。在相同几何尺度下，热量比水分更快影响内部。以半径定义式\eqref{eq:bi}中的Biot数，热Biot数为1.3889，需保留内部温度梯度。

\begin{equation}\label{eq:bi}
\alpha=\frac{k}{\rho_h c_p},\qquad
{\rm Bi}_T=\frac{hR}{k},\qquad
{\rm Bi}_C=\frac{h_mR}{D}.
\end{equation}

初始水分Biot数约为3.2404。将$C=0.15$代入题设函数，得到$D=1.8547\times10^{-11}$平方米每秒，相应水分Biot数约为862.65。这里用函数值说明低含水率时内部阻力的增大，所取$C=0.15$并非问题一在0.50小时的实际状态。计算中仍按各节点状态更新$D$。

表面含水率在短时间内下降较快，表面状态与环境之间仍保留有限差值。Robin条件通过题给传递系数确定交换速率，描述表面温度和含水率对环境变化的响应。

\section{问题二：变物性热湿耦合过程}
\subsection{前三小时的结果与反馈}
问题二的六个指定时刻结果见表\ref{tab:paper3}和表\ref{tab:paper4}。0.50小时，中心、表面温度分别为32.1892、35.4130摄氏度，均低于问题一；表面含水率为1.6486，高于问题一的1.5102。这一差异来自热容量、导热率和扩散函数的共同改变，传热及传质系数保持不变。

\begin{table}[htbp]
\centering\small
\caption{问题二指定时刻的温度（对应题目表3，单位：摄氏度）}\label{tab:paper3}
\begin{tabular}{rrrrrr}
\toprule
时间/h & 0 cm & 0.5 cm & 1.0 cm & 1.5 cm & 2.0 cm \\
\midrule
0.5000 & 32.1892 & 32.3821 & 32.9659 & 33.9608 & 35.4130 \\
1 & 40.3816 & 40.5536 & 41.0605 & 41.8782 & 42.9976 \\
1.5000 & 45.8468 & 45.9348 & 46.1930 & 46.6051 & 47.1401 \\
2 & 48.4502 & 48.4881 & 48.5984 & 48.7736 & 49.0033 \\
2.5000 & 49.4670 & 49.4792 & 49.5136 & 49.5653 & 49.6609 \\
3 & 49.8495 & 49.8553 & 49.8746 & 49.9101 & 49.9664 \\
\bottomrule
\end{tabular}
\end{table}


\begin{table}[htbp]
\centering\small
\caption{问题二指定时刻的干基含水率（对应题目表4，单位：kg/kg）}\label{tab:paper4}
\begin{tabular}{rrrrrr}
\toprule
时间/h & 0~cm & 0.5~cm & 1.0~cm & 1.5~cm & 2.0~cm \\
\midrule
0.5000 & 2.5499 & 2.5489 & 2.5256 & 2.3257 & 1.6486 \\
1.0000 & 2.5257 & 2.4948 & 2.3578 & 2.0230 & 1.4711 \\
1.5000 & 2.3861 & 2.3256 & 2.1344 & 1.8020 & 1.3476 \\
2.0000 & 2.1709 & 2.1084 & 1.9236 & 1.6259 & 1.2311 \\
2.5000 & 1.9566 & 1.9006 & 1.7360 & 1.4720 & 1.1166 \\
3.0000 & 1.7662 & 1.7165 & 1.5702 & 1.3333 & 1.0081 \\
\bottomrule
\end{tabular}
\end{table}


3小时，中心、表面温度分别为49.8495、49.9664摄氏度，径向温差为0.1169摄氏度；含水率分别为1.7662、1.0081，相差0.7581。温度已接近均匀，水分仍有明显梯度。

图\ref{fig:q2-fields}在共同的时间与半径坐标上展示前三小时温湿场。温度由外向内升高，含水率由表面向内降低。升温逐渐接近完成时，水分迁移仍在继续，后续烘干主要受内部扩散及表层阻力影响。

\begin{figure}[htbp]\centering
\includegraphics[width=\linewidth]{figures/result_q2_fields.pdf}
\caption{问题二前三小时的温度与含水率场，按60秒间隔和101个径向重建点展示。}\label{fig:q2-fields}
\end{figure}

\subsection{扩散函数沿计算轨迹的变化}
扩散系数随温度升高而增大，随含水率降低而减小。图\ref{fig:q2-dtrajectory}将求得的温湿场代入题设函数，展示两种作用共同形成的局部扩散系数变化，色标采用对数形式。

\begin{figure}[htbp]\centering
\includegraphics[width=\linewidth]{figures/process_q2_diffusion.pdf}
\caption{问题二前三小时的局部扩散系数，由温湿解代入题设函数得到，色标采用对数形式。}\label{fig:q2-dtrajectory}
\end{figure}

在50摄氏度下，$C=2.55$时的扩散系数约为$1.3471\times10^{-8}$平方米每秒；$C=0.15$时为$8.0012\times10^{-10}$，$C=0.05$时降至$1.9833\times10^{-12}$。即使表层已升温，含水率降低仍可使其成为低扩散区，这有助于解释后期脱水减慢。

含水率降低时，有效体积热容量和导热率均下降。前者使同等净热流引起的温度变化加快，后者减小内部导热能力；两种作用由联立方程共同确定。问题三延续该初边值问题，计算至全域含水率阈值。

问题二、三采用同一积分解。问题二从连续解中逐秒采样，问题三按60秒间隔输出，并由事件函数确定终点。两者的区别在于结果呈现方式。

\section{问题三：固定半径的全域达标时间}
\subsection{事件判据与三种含水指标}
“药材内部各处均低于0.15”对应最大值条件。式\eqref{eq:metrics}分别定义径向最大含水率、体积权均值和表面含水率。计算在全部节点上取最大值，保形径向重建保持节点间的单调性。在已检查时刻，最大值位于中心附近；终点算法仍使用全部节点的最大值。

\begin{equation}\label{eq:metrics}
C_{\max}(t)=\max_{0\leq r\leq R(t)}C(r,t),\qquad
\overline C(t)=\frac{2}{R(t)^2}\int_0^{R(t)}rC(r,t)\,dr,\qquad
C_s(t)=C(R(t),t).
\end{equation}

先定位$C_{\max}-0.15$由正到负的事件，再按式\eqref{eq:event}确定严格小于阈值的整秒输出时刻。由于题目采用严格不等式，事件恰在整数秒发生时仍取下一秒，并用未舍入的全部节点状态检查该时刻。

\begin{equation}\label{eq:event}
C_{\max}(t_*)=0.15\ \text{且向下穿越},\qquad
t_{\rm int}=\lfloor t_*\rfloor+1,\qquad
C_{\max}(t_{\rm int})<0.15.
\end{equation}

图\ref{fig:q3-threshold}比较三种含水指标。按60秒保存轨迹线性插值，表面值约在12.5521小时达到0.15，均值约在36.0251小时达到0.15；最大值的事件时间为57.4728小时。前两个时间是轨迹插值近似值，最大值时间由事件求根得到。采用最大值判据能覆盖内部较慢脱水的位置。

\begin{figure}[htbp]\centering
\includegraphics[width=\linewidth]{figures/process_q3_threshold.pdf}
\caption{问题三的最大、平均及表面含水率，按60秒间隔并补入终点；点线表示阈值0.15。}\label{fig:q3-threshold}
\end{figure}

\subsection{终点数值及每六小时结果}
最终网格下，问题三的越阈时间为57.4728小时。终点表面含水率约为0.0527，体积权均值约为0.1243。表\ref{tab:paper5}列出每6小时及严格终点的指定位置结果。

\begin{table}[htbp]
\centering\small
\caption{问题三每6 h及严格终点的干基含水率（对应题目表5）}\label{tab:paper5}
\begin{tabular}{rrrrrr}
\toprule
时间/h & 0~cm & 0.5~cm & 1.0~cm & 1.5~cm & 2.0~cm \\
\midrule
6 & 1.0170 & 0.9881 & 0.9015 & 0.7550 & 0.5328 \\
12 & 0.4561 & 0.4432 & 0.4031 & 0.3297 & 0.1642 \\
18 & 0.2988 & 0.2915 & 0.2687 & 0.2253 & 0.0845 \\
24 & 0.2378 & 0.2327 & 0.2167 & 0.1855 & 0.0665 \\
30 & 0.2058 & 0.2018 & 0.1892 & 0.1641 & 0.0598 \\
36 & 0.1859 & 0.1825 & 0.1718 & 0.1504 & 0.0566 \\
42 & 0.1721 & 0.1691 & 0.1597 & 0.1407 & 0.0549 \\
48 & 0.1618 & 0.1592 & 0.1507 & 0.1334 & 0.0537 \\
54 & 0.1539 & 0.1514 & 0.1436 & 0.1277 & 0.0530 \\
57.4728 & 0.1500 & 0.1477 & 0.1402 & 0.1249 & 0.0527 \\
\bottomrule
\end{tabular}
\par\smallskip\begin{minipage}{0.93\linewidth}\footnotesize 数值单位为 kg/kg。末行使用严格整秒终点，时间转为小时后显示四位小数；中心显示 0.1500 是舍入结果，严格不等式按未舍入值检查。\end{minipage}
\end{table}


表\ref{tab:paper5}显示，中心含水率由6小时的1.0170降至12小时的0.4561，48小时仍为0.1618，54小时为0.1539。接近阈值时，绝对剩余差值虽小，仍需较长时间继续脱水。这与低含水率区扩散系数迅速下降的特征一致。

图\ref{fig:q3-terminal}给出终点径向剖面。中心附近接近阈值，表面接近且略高于有效边界值0.05。表面与环境间的正含水率差维持向外传质，使内部水分继续排出。

\begin{figure}[htbp]\centering
\includegraphics[width=\linewidth]{figures/result_q3_terminal.pdf}
\caption{问题三终点的径向含水率分布，中心接近阈值，表面含水率较低。}\label{fig:q3-terminal}
\end{figure}

\section{问题四：实测收缩下的移动边界}
\subsection{随体坐标与运动项}
在均匀径向收缩假设下，材料速度为$v(r,t)=r\dot R/R$。令$x=r/R(t)$，对随材料输运的标量$u$定义$U(x,t)=u(R(t)x,t)$，可得式\eqref{eq:ale}。固定空间时间导数中的坐标变化项与材料平流项抵消，随体坐标下的时间导数即为物质导数。

\begin{equation}\label{eq:ale}
\left.\frac{\partial u}{\partial t}\right|_r
=\left.\frac{\partial U}{\partial t}\right|_x
-\frac{x\dot R}{R}U_x,\qquad
v\,u_r=\frac{x\dot R}{R}U_x,\qquad
\frac{{\rm D}u}{{\rm D}t}=U_t.
\end{equation}

移动边界干燥中可用ALE方法描述网格和材料运动\cite{seyedabadi2017ale}。这里采用网格随材料均匀收缩的情形，固定$x$对应同一材料位置。

随体坐标下的方程和表面条件见式\eqref{eq:moving}，其中时间导数均在固定$x$下计算。沿用有限体积离散，在各时刻更新半径和问题四物性。内部扩散含$R(t)^{-2}$因子，归一化表面交换含$R(t)^{-1}$因子，二者共同决定收缩后的传递速度。

\begin{equation}\label{eq:moving}
\begin{aligned}
A(C)\theta_t&=\frac{1}{R(t)^2x}
                \partial_x\!\left(xk(C)\theta_x\right),\\
C_t&=\frac{1}{R(t)^2x}
                \partial_x\!\left(xD(\theta,C)C_x\right),\\
-k_s\theta_x(1,t)&=R(t)h(\theta_s-\theta_\infty),\\
-D_sC_x(1,t)&=R(t)h_m(C_s-C_e).
\end{aligned}
\end{equation}

\subsection{干基变量为何不直接增加压缩项}
干基含水率为水质量与干物质质量之比。均匀收缩使单位体积内两种质量同时变化，压缩本身保持该比值。设干固体密度$\rho_d(t)$在截面上均匀，长度不变且干质量守恒，则$\rho_dR^2$为常数。取相对干固体的扩散通量为$-\rho_dD\nabla C$，由水分与干固体连续性方程得到式\eqref{eq:drybasis}。

\begin{equation}\label{eq:drybasis}
\begin{aligned}
\frac{{\rm D}\rho_d}{{\rm D}t}+\rho_d\nabla\!\cdot v&=0,\\
\rho_d\frac{{\rm D}C}{{\rm D}t}
&=\nabla\!\cdot(\rho_dD\nabla C),\\
\rho_d(t)=\rho_{d0}\left(\frac{R_0}{R(t)}\right)^2\quad&\Longrightarrow\quad
\frac{{\rm D}C}{{\rm D}t}=\nabla\!\cdot(D\nabla C).
\end{aligned}
\end{equation}

空间均匀的$\rho_d(t)$可从散度中提出并约去，得到式\eqref{eq:moving}中的水分方程，无需附加$-C\nabla\!\cdot v$项。初始密度$\rho_{d0}>0$也在归一化方程中约去，计算无需另赋数值。若改用单位当前体积的水浓度，压缩项的形式将随变量定义改变；相关移动边界研究采用过这类体积变量\cite{adrover2020pears}。

题给$\rho_h=760+90C$按表\ref{tab:closure}用于有效热容量，水分收支则相对于不变干质量归一化。两种参数分别承担能量闭合和质量守恒的作用。

\subsection{含水率场、域外位置与严格终点}
图\ref{fig:q4-moving}以物理半径为纵轴，黑线表示插值半径$R(t)$，域外留白。固定物理半径对应水平线，而同一材料点随$r=xR(t)$运动，应沿该轨迹读取其状态变化。

\begin{figure}[htbp]\centering
\includegraphics[width=\linewidth]{figures/process_q4_moving.pdf}
\caption{问题四的移动域含水率场，按60秒间隔并补入终点；黑线为当前表面，域外留白。}\label{fig:q4-moving}
\end{figure}

采用与问题三相同的最大含水率判据，问题四的越阈时间为51.0908小时。终点表面含水率约为0.0526，体积权均值约为0.1184，半径为1.2000厘米。终点位于半径记录的72小时范围内。

表\ref{tab:paper6}列出每6小时及严格终点的指定位置含水率。在这些表列时刻，0、0.5、1.0厘米均位于药材内，1.5厘米均在域外。表面列对应实际$R(t)$处状态，随半径轨迹变化。工作簿4中的20099个域外位置按定义留空。

\begin{table}[htbp]
\centering\small
\caption{问题四每6小时及严格终点的干基含水率（对应题目表6）}\label{tab:paper6}
\begin{tabular}{rrrrrr}
\toprule
时间/h & 0 cm & 0.5 cm & 1.0 cm & 1.5 cm & 移动表面 \\
\midrule
6 & 1.7196 & 1.5377 & 1.0226 & --- & 0.4207 \\
12 & 0.7377 & 0.6546 & 0.4083 & --- & 0.1670 \\
18 & 0.4087 & 0.3686 & 0.2397 & --- & 0.0892 \\
24 & 0.2851 & 0.2611 & 0.1790 & --- & 0.0673 \\
30 & 0.2263 & 0.2094 & 0.1493 & --- & 0.0595 \\
36 & 0.1928 & 0.1797 & 0.1317 & --- & 0.0560 \\
42 & 0.1712 & 0.1604 & 0.1200 & --- & 0.0541 \\
48 & 0.1561 & 0.1468 & 0.1116 & --- & 0.0530 \\
51.0908 & 0.1500 & 0.1413 & 0.1081 & --- & 0.0526 \\
\bottomrule
\end{tabular}
\par\smallskip\begin{minipage}{0.93\linewidth}\footnotesize “---”表示该时刻位置已在药材外，结果工作簿对应单元格为空；数值单位为 kg/kg。末行使用严格整秒终点，时间转为小时后显示四位小数；中心显示 0.1500 是舍入结果，严格不等式按未舍入值检查。\end{minipage}
\end{table}



问题四在6小时的中心含水率为1.7196，高于问题三的1.0170；42小时为0.1712，略低于问题三的0.1721。按60秒最大含水率轨迹插值，两曲线约在41.2小时交叉。物性变化与扩散距离缩短共同影响全过程，问题四在早期较湿，最终却较早达到阈值。
